\documentclass[conference]{IEEEtran}
\usepackage{cite}
\usepackage{amsmath,amssymb,amsfonts}
\usepackage{algorithmic}
\usepackage{graphicx}
\usepackage{textcomp}
\usepackage{xcolor}
\usepackage{booktabs}
\def\BibTeX{{\rm B\kern-.05em{\sc i\kern-.025em b}\kern-.08em
    T\kern-.1667em\lower.7ex\hbox{E}\kern-.125emX}}
\begin{document}

\title{Early Fault Detection and Risk-Based Maintenance Prioritization for Wind Turbines}

\author{
\IEEEauthorblockN{Ziyad Muradov}
\IEEEauthorblockA{\textit{AI Academy} \\
\textit{National Artificial Intelligence Center}\\
Baku, Azerbaijan}
\and
\IEEEauthorblockN{Emin Huseynli}
\IEEEauthorblockA{\textit{AI Academy} \\
\textit{National Artificial Intelligence Center}\\
Baku, Azerbaijan}
\and
\IEEEauthorblockN{Kamal Soltanaliyev}
\IEEEauthorblockA{\textit{AI Academy} \\
\textit{National Artificial Intelligence Center}\\
Baku, Azerbaijan}
\and
\IEEEauthorblockN{Ismayil Yusifli}
\IEEEauthorblockA{\textit{AI Academy} \\
\textit{National Artificial Intelligence Center}\\
Baku, Azerbaijan}
}

\maketitle

\begin{abstract}
% TODO (whoever finalizes, likely Kamal per the report-driving role):
% write the final 150-200 word abstract once Results and Discussion are
% complete. Draft below for reference/replacement.
Reactive wind turbine maintenance, triggered only after a fault is
formally logged, forgoes the opportunity for early intervention. This
work investigates whether a recurrent neural network trained on SCADA
sensor windows can detect an approaching fault early enough, and with
sufficient calibrated confidence, to support risk-based maintenance
prioritization across a multi-farm fleet, and whether such a model
generalizes to turbines and farms it was not trained on. We frame the
task as per-timestep risk regression using a Gated Recurrent Unit over a
farm-agnostic physical feature representation, evaluate it against
classical baselines and across multiple farms, and report an honest,
evidenced account of where the approach succeeds and where domain shift
between farms limits it.

GitHub repository: \url{https://github.com/thedarkstringg/team-Sigmoid}
(tag \texttt{v1.0-final}).
\end{abstract}

\begin{IEEEkeywords}
wind turbine, fault detection, predictive maintenance, recurrent neural
network, domain generalization, SCADA
\end{IEEEkeywords}

\section{Introduction}
% TODO (owner: whole team, or assign one person - not yet claimed as of
% this commit). Motivate the problem, state the research question, and
% preview the paper's structure and headline finding.

\section{Related Work}
% TODO (owner: Ismayil, per team split). SCADA fault-detection literature
% review, including comparison against the CARE-to-Compare ~48h benchmark
% referenced in \cite{gueck2024}.

\section{Method}
\label{sec:method}

\subsection{Task Framing}

We frame early fault detection as a per-timestep risk regression problem
rather than a single window-level classification. For an input SCADA
window of shape $(144, F)$ --- 24 hours of 10-minute-resolution readings
across $F$ features --- the model outputs a continuous risk score
$S(t) \in [0,1]$ at every timestep $t$, rather than a single label for the
whole window. This framing is required for the lead-time evaluation
(Section~\ref{sec:results}): computing lead time as the interval between a
predicted alarm and the actual fault requires identifying the specific
timestep at which predicted risk crosses an alarm threshold $\tau$, which a
single window-level prediction cannot provide.

\subsection{Feature Representation}

Raw SCADA sensor channels differ completely in count and naming across
farms (54--238 average-statistic sensors, depending on farm), which makes
a model trained on one farm's raw sensors architecturally incompatible
with another farm's data. To support the cross-farm generalization study
(Section~\ref{sec:results}), we instead map each farm's raw sensors onto a
shared set of physically-meaningful, farm-agnostic features, computed
identically regardless of which farm the data originates from. The
canonical representation used for all cross-farm comparisons consists of
ten features: wind speed, blade pitch angle, yaw misalignment, gearbox oil
temperature rise above ambient, gearbox bearing hotspot over oil
temperature, generator-to-rotor speed ratio, grid power factor, and three
grid electrical-imbalance/deviation terms (current imbalance, voltage
imbalance, frequency deviation). An eleventh feature, power residual
(actual power minus a wind-speed-conditioned expected power curve), was
evaluated as a separate ablation but excluded from the canonical
representation, since fitting its expected-power curve requires
target-domain data and would violate the zero-shot constraint of the
cross-farm study.

Each feature is standardized (zero mean, unit variance) using statistics
fit exclusively on the source farm's training split; the same frozen
scaler is applied unchanged to validation, test, and any external target
farm's data, so that no target-domain information influences the
input representation.

\subsection{Sequence Model}

We use a Gated Recurrent Unit (GRU) network rather than a Long Short-Term
Memory (LSTM) network or a Transformer. The dataset is modest in size
relative to typical deep learning corpora (on the order of $10^3$--$10^4$
training windows depending on farm); a GRU has fewer learned parameters
than an LSTM of equal hidden size, since it merges the forget and input
gates into a single update gate and omits the separate cell state,
reducing overfitting risk on limited data. Transformers, by contrast,
generally require substantially more training data than recurrent
architectures to outperform them, since they lack a built-in sequential
inductive bias; we judged this an unfavorable trade-off at this dataset
scale.

The final architecture is a 2-layer GRU with hidden size 32 and dropout
0.3 between layers, followed by a linear projection applied identically at
every timestep, mapping the hidden state to a single logit:
\begin{align}
h_t &= \mathrm{GRU}(x_t, h_{t-1}) \\
z_t &= W h_t + b \\
S(t) &= \sigma(z_t)
\end{align}
where $\sigma$ denotes the sigmoid function. The model returns raw logits
$z_t$ during training; $\sigma(z_t)$ is applied separately at inference and
evaluation time. Architecture depth (2 vs. 3 layers) and a bidirectional
variant were both evaluated as ablations (Section~\ref{sec:results}) and
did not improve validation performance sufficiently to justify their added
complexity or, in the bidirectional case, the loss of real-time causal
deployability.

\subsection{Loss Function and Class Imbalance}

We use binary cross-entropy computed jointly with the sigmoid
(\texttt{BCEWithLogitsLoss}) rather than composing a separate sigmoid
layer with plain binary cross-entropy. The combined formulation avoids
computing $\log$ of values arbitrarily close to 0 or 1 directly, which is
numerically unstable in floating-point arithmetic and can produce vanishing
or exploding gradients as the model becomes confident.

Fault-approaching timesteps are rare relative to normal operation
(positive rate ranging from approximately 2\% to 6\% depending on farm and
split). We address this with a positive-class weight, computed as the
ratio of negative to positive timesteps within the training split only,
to avoid leaking split-level statistics into the loss function.

\subsection{Masking}

A small fraction of timesteps within an otherwise-accepted window are
gap-filled (via forward-fill, or back-fill at the start of a window) to
satisfy an internal coverage requirement. Gap-filled timesteps are marked
with a binary observation mask and excluded from the loss calculation, by
multiplying the per-timestep loss by the mask before averaging:
\begin{equation}
\mathcal{L} = \frac{\sum_t \mathcal{L}_t \cdot m_t}{\sum_t m_t}
\end{equation}
where $m_t \in \{0, 1\}$. The same mask is applied when computing
evaluation metrics, so reported results reflect performance on real
observations only. Certain physical features (e.g., the generator-to-rotor
speed ratio) can produce extreme values even when correctly marked
invalid; as a defensive stopgap, standardized inputs are additionally
clipped to $[-10, 10]$ before being passed to the model, since a single
extreme value elsewhere in a window can otherwise distort the recurrent
hidden state at other, valid timesteps in the same window.

\subsection{Compute Budget}

Training uses mixed-precision arithmetic and gradient accumulation to
control memory footprint, and a bounded 24-hour (144-timestep) input
window rather than a longer history, in line with the project's compute
constraints. In practice, the combination of a small dataset and a
compact 2-layer, 32-hidden-unit GRU resulted in peak GPU memory usage and
training time well within the allotted budget (on the order of seconds
per epoch; see the deployment note in Section~\ref{sec:results} for
measured latency and memory figures).

\subsection{Hyperparameter Selection}

We evaluated six configurations, varying learning rate
($1\times10^{-4}$--$5\times10^{-4}$), dropout (0.3--0.5), feature count (10
vs.\ 11, i.e.\ with the power-residual ablation), and architecture
(unidirectional vs.\ bidirectional). The final configuration was selected
by lowest validation loss, without reference to test-set performance, to
avoid selecting a configuration that happens to fit the test split rather
than generalizing. All six configurations showed a consistent pattern:
strong validation performance (ROC-AUC 0.57--0.67 across configurations)
alongside markedly weaker held-out test performance (ROC-AUC 0.40--0.48),
a finding discussed in detail in Section~\ref{sec:discussion}.

\subsection{Reproducibility}

All experiments use a fixed random seed applied to Python's, NumPy's, and
PyTorch's random number generators. Model and optimizer state are
checkpointed after every epoch, with the best-validation-loss checkpoint
retained separately from the most recent one, to tolerate the shared
compute environment's training window ending mid-run without losing
progress or requiring a full restart.

\section{Experimental Setup}
\label{sec:experimental_setup}
% TODO (owner: Ismayil, per team split). Dataset description (CARE-to-
% Compare, licence, size per farm), splits, baseline model details
% (Logistic Regression / HistGradientBoosting), class-imbalance handling.

\section{Results}
\label{sec:results}
% TODO (owner: Emin, per team split). Quantitative result table (val+test,
% all farms), lead-time metric, per-turbine/per-farm breakdown, cost-
% sensitive metric, calibration analysis, deployment note (latency, model
% size, memory), Dockerfile reference. Required plots per the brief:
% reliability diagram, lead-time distribution, per-turbine breakdown chart.

\section{Discussion and Limitations}
\label{sec:discussion}
% TODO (owner: Kamal, per team split). Cross-farm generalization findings,
% distribution-shift analysis (Tasks 16-18), turbine-model shift, label-
% derivation differences across farms, compute/scoping limitations tying
% back to Section~\ref{sec:method}.

\section{Conclusion}
% TODO (owner: not yet assigned as of this commit). Summarize the
% research question, headline finding, and honest takeaway.

\section*{Tools and Acknowledgements}
% TODO: each team member should add a brief, honest note on any
Claude was used to independently verify that src/ implementation and tests/ pass. The code in src/ itself was written independently (not generated by AI). Per the course’s oral verification policy (Week 7), we are responsible for being able to
explain and reproduce every derivation and code reference in
this report without AI assistance.


\begin{thebibliography}{00}
\bibitem{gueck2024} C. G\"{u}ck and C. Roelofs, ``Wind Turbine SCADA Data
For Early Fault Detection (v1.0),'' Zenodo, 2024. Available:
\url{https://doi.org/10.5281/zenodo.10958775}
% TODO (owner: Ismayil): add remaining citations for the Related Work
% literature review.
\end{thebibliography}

\end{document}